\subsection{SPECIAL CRITERIA FOR EQUICONTINUITY}

It is clear that every subset of an equicontinuous (resp. uniformly equicontinuous) set is equicontinuous (resp. uniformly equicontinuous). Again, if X is a topological (resp. uniform) space and Y is a uniform space, every finite union of equicontinuous (resp. uniformly equicontinuous) subsets of \(\mathcal{F}(X;Y)\) is equicontinuous (resp. uniformly equicontinuous).

Let \(X, X'\) be two topological (resp. uniform) spaces, let \(Y, Y'\) be two uniform spaces, let \(f: X \to X'\) be a continuous (resp. uniformly continuous) mapping and let \(g: Y \to Y'\) be a uniformly continuous mapping.

It follows immediately from the definitions that the mapping \(u \rightarrow g \circ u \circ f\) of \(\mathcal{F}(X; Y)\) into \(\mathcal{F}(X'; Y')\) transforms equicontinuous (resp. uniformly equicontinuous) sets into equicontinuous (resp. uniformly equicontinuous) sets.

PROPOSITION 3. Let X be a topological (resp. uniform) space, let \((\mathbf{Y}_{t})_{t \in I}\) be a family of uniform spaces, let Y be a set, and for each \(t \in I\) , let \(f_{t}\) be a mapping of Y into \(\mathbf{Y}_{t}\) . Let Y be endowed with the coarsest uniformity for which all the \(f_{t}\) are uniformly continuous. For a subset H of \(\mathcal{F}(\mathbf{X}; \mathbf{Y})\) to be equicontinuous (resp. uniformly equicontinuous) it is necessary and sufficient that, for each \(t \in I\) , the image of H under the mapping \(u \to f_{t} \circ u\) be an equicontinuous (resp. uniformly equicontinuous) subset of \(\mathcal{F}(\mathbf{X}; \mathbf{Y}_{t})\) .

This is an immediate consequence of Definitions 1 and 2 and the definition of the entourages of Y.

PROPOSITION 4. Let X, Y be two uniform spaces and let H be a set of uniformly continuous mappings of X into Y. Let \(\hat{X}, \hat{Y}\) be the Hausdorff completions of X, Y respectively, and let \(\tilde{H}\) denote the set of mappings \(\hat{u}: \hat{X} \to \hat{Y}\) as \(u\) runs through H (Chapter II, § 3, no. 7, Proposition 15). Then H is uniformly equicontinuous if and only if \(\tilde{H}\) is uniformly equicontinuous.

We recall that the diagram

\[
\begin{array}{c} \mathbf {X} \xrightarrow {u} \mathbf {Y} \\ i \downarrow \quad \quad \quad \quad \quad \quad \quad \quad \quad \quad \quad \quad \quad j \\ \hat {\mathbf {X}} \xrightarrow {\hat {u}} \hat {\mathbf {Y}} \end{array}\tag{1}
\]

is commutative, where \(i\) and \(j\) are the canonical mappings; moreover, the uniformity of \(\mathbf{X}\) (resp. \(\mathbf{Y}\) ) is the inverse image under \(i\) (resp. \(j\) ) of that of \(\hat{\mathbf{X}}\) (resp. \(\hat{\mathbf{Y}}\) ). Hence \(\mathbf{H}\) is uniformly equicontinuous if and only if its image under the mapping \(u \to j \circ u\) is uniformly equicontinuous (Proposition 3), and we may already restrict ourselves to the case where \(\mathbf{Y}\) is Hausdorff and complete; moreover, if \(\tilde{\mathbf{H}}\) is uniformly equicontinuous then so is \(\mathbf{H}\) , because it is the image of \(\tilde{\mathbf{H}}\) under the mapping \(\hat{u} \to \hat{u} \circ i\) ; thus it remains to prove the converse when \(\mathbf{Y} = \hat{\mathbf{Y}}\) . Let \(\mathbf{V}\) be a closed entourage of \(\mathbf{Y}\) ; by hypothesis there is an entourage \(\mathbf{U}\) of \(\mathbf{X}\) such that the relations \((x, x') \in \mathbf{U}\) and \(u \in \mathbf{H}\) imply that \((u(x), u(x')) \in \mathbf{V}\) . Now, if \(\mathbf{U}'\) is the image of \(\mathbf{U}\) under \(i \times i\) , the closure \(\overline{\mathbf{U}'}\) of \(\mathbf{U}'\) in \(\hat{\mathbf{X}} \times \hat{\mathbf{X}}\) is an entourage of \(\hat{\mathbf{X}}\) (Chapter II, § 3, no. 7, Proposition 12); the hypothesis implies that, whenever \((z, z') \in \mathbf{U}'\) and \(u \in \mathbf{H}\) , we have \((\hat{u}(z), \hat{u}(z')) \in \mathbf{V}\) . Since \(\mathbf{V}\) is closed and \(\hat{u}\) is continuous, we have also \((\hat{u}(t), \hat{u}(t')) \in \mathbf{V}\) for all \((t, t') \in \overline{\mathbf{U}'}\) and all \(u \in \mathbf{H}\) ; this completes the proof.

PROPOSITION 5. Let G, G' be two topological groups endowed with their left uniformities, and let H be a set of homomorphisms of G into G'. Then the following conditions are equivalent:

\begin{enumerate}
\def\labelenumi{\alph{enumi})}
\item
  H is equicontinuous at the identity element e of G,
\item
  H is equicontinuous,
\item
  H is uniformly equicontinuous.
\end{enumerate}

It is enough to show that a) implies c). Let \(V'\) be a neighbourhood of the identity element \(e'\) of \(G'\) ; then, by hypothesis, there is a neighbourhood V of e in G such that \(u(V) \subset V'\) for all \(u \in H\) ; since the elements of H are homomorphisms, the relation \(x^{-1}y \in V\) implies that we have

\[
(u (x)) ^ {- 1} u (y) = u (x ^ {- 1} y) \in \mathbf {V} ^ {\prime}.
\]

In view of the definition of the entourages of the left uniformities of G and G' (Chapter III, § 3, no. 1), the result follows.

